\section*{1. Matrices and column combinations}

\notetarget{note:01:matrices}A real \textbf{matrix} \(A=(a_{ij})\in\R^{m\times n}\) has \(m\) rows and
\(n\) columns; \(a_{ij}\) is its entry in row \(i\), column \(j\).
Matrices are equal when their sizes and corresponding entries agree.
A vector in \(\R^n\) is a column:
\[
x=\begin{pmatrix}x_1\\\vdots\\x_n\end{pmatrix}.
\]
For matrices of the same size and real \(c\), define
\[
(A+B)_{ij}=a_{ij}+b_{ij},\qquad(cA)_{ij}=ca_{ij}.
\]
The zero matrix, denoted by \(0\), has every entry zero; its size is determined
by the equation in which it occurs. Set \(-A=(-1)A\) and \(A-B=A+(-B)\).
\notetarget{note:01:coordinate-laws}These operations satisfy
\begin{align*}
A+B&=B+A,&(A+B)+C&=A+(B+C),\\
c(A+B)&=cA+cB,&(c+d)A&=cA+dA,\\
c(dA)&=(cd)A,&1A&=A,\\
A+0&=A,&A+(-A)&=0.
\end{align*}
\notetarget{note:01:coordinate-laws-proof}Each identity holds entry by entry by the corresponding real-number law;
for example, \(c(a_{ij}+b_{ij})=ca_{ij}+cb_{ij}\).

\notetarget{note:01:linear-combinations}A \textbf{linear combination} of same-sized matrices \(A_1,\ldots,A_k\)
is \(c_1A_1+\cdots+c_kA_k\), where the coefficients are real.
The empty combination is defined to be zero. A vector is a special case.
\notetarget{note:01:standard-columns}Let \(e_j\in\R^n\) have entry \(1\) in position \(j\) and zeros elsewhere.
\notetarget{note:01:standard-coordinates-proof}Comparing coordinate \(i\) gives
\[
x=x_1e_1+\cdots+x_ne_n.
\]

\subsection*{The product \(Ax\)}
\notetarget{note:01:ax}Write \(A=[\,a_1\ \cdots\ a_n\,]\), with \(a_j\in\R^m\).
For \(x\in\R^n\), define
\[
\boxed{Ax=x_1a_1+\cdots+x_na_n.}
\]
Its \(i\)th coordinate is \(\sum_{j=1}^n a_{ij}x_j\). Consequently
\[
Ax=b\quad\Longleftrightarrow\quad
\sum_{j=1}^n a_{ij}x_j=b_i\quad(1\leqslant i\leqslant m).
\]
\notetarget{note:01:linear-system}This is a \textbf{linear system}: \(A\) is the coefficient matrix,
\(b\in\R^m\) is its right-hand side, and \(x\) is the unknown column.
A solution must satisfy every equation. The system is \textbf{consistent}
if a solution exists. For \(b=0\) it is \textbf{homogeneous}; then \(x=0\)
is always a solution. A \textbf{nontrivial} solution is a nonzero column.

\begin{example}[Representing a column]\label{example:01:representation}\notetarget{note:01:representation}
Let
\[
A=\begin{pmatrix}2&-4&0&0\\-1&2&0&0\\0&0&1&2\end{pmatrix},\qquad
b=\begin{pmatrix}4\\-2\\-3\end{pmatrix},\qquad
c=\begin{pmatrix}1\\0\\0\end{pmatrix}.
\]
Determine whether \(Ax=b\) and \(Ax=c\) have solutions.
\notetarget{note:01:representation-solution}The equality
\[
2\begin{pmatrix}2\\-1\\0\end{pmatrix}
+0\begin{pmatrix}-4\\2\\0\end{pmatrix}
-3\begin{pmatrix}0\\0\\1\end{pmatrix}
+0\begin{pmatrix}0\\0\\2\end{pmatrix}=b
\]
shows that \(Ax=b\) has a solution.
For every \(x\), coordinate 1 of \(Ax\) is minus twice coordinate 2:
\[
2x_1-4x_2=-2(-x_1+2x_2).
\]
The column \(c\) fails this condition, since \(1\ne-2(0)\).
Thus \(Ax=c\) has no solution.
\end{example}

\begin{proposition}[Linearity of multiplication by a matrix]\label{prop:01:linearity}\notetarget{note:01:linearity}
For \(A\in\R^{m\times n}\), \(x,y\in\R^n\), and \(s,t\in\R\),
\[
A(sx+ty)=sAx+tAy.
\]
\end{proposition}
\begin{proof}\notetarget{note:01:linearity-proof}
By the column definition and the coordinate laws,
\[
A(sx+ty)=\sum_j(sx_j+ty_j)a_j
=s\sum_jx_ja_j+t\sum_jy_ja_j=sAx+tAy.
\]
\end{proof}
\notetarget{note:01:finite-linearity}The same calculation gives \(A\sum_{j=1}^k c_jx_j=\sum_{j=1}^k c_jAx_j\)
for any finite list: apply the two-term identity successively, with the
empty case \(A0=0\).

\section*{2. Matrix multiplication and transpose}

\notetarget{note:01:matrix-product}For \(A\in\R^{m\times n}\) and \(B=[\,b_1\ \cdots\ b_p\,]\in\R^{n\times p}\),
define
\[
AB=[\,Ab_1\ \cdots\ Ab_p\,]\in\R^{m\times p}.
\]
Column \(j\) of \(B\) supplies the coefficients for column \(j\) of \(AB\).
In entries,
\[
\boxed{(AB)_{ij}=\sum_{k=1}^n a_{ik}b_{kj}.}
\]
The number of columns of \(A\) must equal the number of rows of \(B\).

\begin{example}[A rectangular product]\label{example:01:product}\notetarget{note:01:rectangular-product}
\notetarget{note:01:product-entries}For
\[
A=\begin{pmatrix}2&-3&0\\1&4&2\end{pmatrix},\qquad
B=\begin{pmatrix}3&4\\-2&0\\1&2\end{pmatrix},
\]
\[
AB=\begin{pmatrix}
2(3)+(-3)(-2)+0(1)&2(4)+(-3)(0)+0(2)\\
1(3)+4(-2)+2(1)&1(4)+4(0)+2(2)
\end{pmatrix}
=\begin{pmatrix}12&8\\-3&8\end{pmatrix}.
\]
\notetarget{note:01:product-columns}The two columns can also be obtained directly as combinations of the
columns of \(A\):
\[
3\begin{pmatrix}2\\1\end{pmatrix}
-2\begin{pmatrix}-3\\4\end{pmatrix}
+\begin{pmatrix}0\\2\end{pmatrix}
=\begin{pmatrix}12\\-3\end{pmatrix},\qquad
4\begin{pmatrix}2\\1\end{pmatrix}
+2\begin{pmatrix}0\\2\end{pmatrix}
=\begin{pmatrix}8\\8\end{pmatrix}.
\]
\end{example}

\begin{proposition}[Product laws]\label{prop:01:product-laws}\notetarget{note:01:product-laws}
For compatible matrix sizes and real \(c\),
\begin{gather*}
A(B+C)=AB+AC,\qquad (A+B)C=AC+BC,\\
(cA)B=A(cB)=c(AB),\qquad (AB)C=A(BC).
\end{gather*}
\end{proposition}
\begin{proof}\notetarget{note:01:product-laws-proof}
The first two identities follow from
\begin{align*}
\bigl(A(B+C)\bigr)_{ij}
 &=\sum_k a_{ik}(b_{kj}+c_{kj})
 =\sum_k a_{ik}b_{kj}+\sum_k a_{ik}c_{kj},\\
\bigl((A+B)C\bigr)_{ij}
 &=\sum_k(a_{ik}+b_{ik})c_{kj}
 =\sum_k a_{ik}c_{kj}+\sum_k b_{ik}c_{kj}.
\end{align*}
\notetarget{note:01:product-scaling-proof}For scalar compatibility,
\[
\sum_k(ca_{ik})b_{kj}=c\sum_k a_{ik}b_{kj}
=\sum_k a_{ik}(cb_{kj}).
\]
\notetarget{note:01:associativity-proof}For associativity take sizes \(m\times n\), \(n\times p\), \(p\times q\).
Both sides have size \(m\times q\), and
\begin{align*}
\bigl((AB)C\bigr)_{ij}
 &=\sum_{k=1}^p\left(\sum_{\ell=1}^n a_{i\ell}b_{\ell k}\right)c_{kj}\\
 &=\sum_{\ell=1}^n a_{i\ell}
       \left(\sum_{k=1}^p b_{\ell k}c_{kj}\right)
 =\bigl(A(BC)\bigr)_{ij}.
\end{align*}
The sums are finite, so their terms may be regrouped. Equality of all
entries proves the claims.
\end{proof}
\notetarget{note:01:product-action-zero}In particular, \((AB)x=A(Bx)\): the factor nearest \(x\) acts first.
The product of a matrix with a compatible zero matrix is zero, since
each entry is a sum of zero terms.

\subsection*{Identity and powers}
\notetarget{note:01:identity}The \textbf{identity matrix} \(I_n=[\,e_1\ \cdots\ e_n\,]\) has ones on
the diagonal and zeros elsewhere. The column definition gives
\(Ae_j=a_j\), so \(AI_n=A\). Also, the \((i,j)\) entry of \(I_mA\)
is \(a_{ij}\), because only the \(i\)th term of its entry sum survives.
Thus
\[
I_mA=A=AI_n\qquad(A\in\R^{m\times n}).
\]
\notetarget{note:01:powers}A matrix is \textbf{square} when its two sizes agree. For square \(A\),
set \(A^0=I\) and \(A^{k+1}=A^kA\) for integers \(k\geqslant0\).
Then
\[
A^rA^s=A^{r+s}\qquad(r,s\geqslant0).
\]
\notetarget{note:01:powers-proof}Indeed, for fixed \(r\) the case \(s=0\) is the identity law. If the
formula holds for \(s\), associativity gives
\(A^rA^{s+1}=A^r(A^sA)=(A^rA^s)A=A^{r+s+1}\).

\begin{example}[Noncommutativity and failure of cancellation]\label{example:01:noncommute}\notetarget{note:01:noncommutativity}
Direct multiplication gives
\[
\begin{pmatrix}1&0\\0&2\end{pmatrix}
\begin{pmatrix}0&1\\1&0\end{pmatrix}
=\begin{pmatrix}0&1\\2&0\end{pmatrix},\qquad
\begin{pmatrix}0&1\\1&0\end{pmatrix}
\begin{pmatrix}1&0\\0&2\end{pmatrix}
=\begin{pmatrix}0&2\\1&0\end{pmatrix}.
\]
Hence \(AB\) need not equal \(BA\), even for square matrices of the same
size. Such matrices \textbf{commute} when \(AB=BA\).

\notetarget{note:01:cancellation}For \(N=\begin{pmatrix}0&1\\0&0\end{pmatrix}\), multiplication gives
\(N^2=0\), although \(N\ne0\). Thus \(NN=N0\) has unequal right factors:
the common nonzero left factor \(N\) cannot be cancelled.
\end{example}

\subsection*{Transpose}
\notetarget{note:01:transpose}The \textbf{transpose} of \(A\in\R^{m\times n}\) is the matrix
\(A^\top\in\R^{n\times m}\) with \((A^\top)_{ij}=a_{ji}\).
\notetarget{note:01:transpose-example}For example,
\[
\begin{pmatrix}2&-3&0\\1&4&2\end{pmatrix}^{\!\top}
=\begin{pmatrix}2&1\\-3&4\\0&2\end{pmatrix}.
\]
A square matrix is \textbf{symmetric} if \(A^\top=A\).
\notetarget{note:01:outer-products}For columns \(x\in\R^m\), \(y\in\R^n\), the \textbf{outer product}
\(xy^\top\) is the \(m\times n\) matrix with entries \(x_i y_j\).
For equal lengths, \(x^\top y=\sum_i x_i y_i\); we identify a
\(1\times1\) matrix with its single entry.

\begin{proposition}[Transpose laws]\label{prop:01:transpose}\notetarget{note:01:transpose-laws}
For compatible sizes,
\[
(A^\top)^\top=A,\quad(A+B)^\top=A^\top+B^\top,\quad
(cA)^\top=cA^\top,\quad (AB)^\top=B^\top A^\top.
\]
\end{proposition}
\begin{proof}\notetarget{note:01:transpose-laws-proof}
The \((i,j)\) entries in the first three identities are, respectively,
\(a_{ij}\), \(a_{ji}+b_{ji}\), and \(ca_{ji}\), on both sides.
\notetarget{note:01:transpose-product-proof}If \(A\) is \(m\times n\) and \(B\) is \(n\times p\), the last two
matrices to compare have size \(p\times m\), and
\[
\bigl((AB)^\top\bigr)_{ij}
=(AB)_{ji}=\sum_{k=1}^n a_{jk}b_{ki}
=\sum_{k=1}^n(B^\top)_{ik}(A^\top)_{kj}
=(B^\top A^\top)_{ij}.
\]
\end{proof}

\section*{3. Elimination and all solutions}

\notetarget{note:01:row-operations}The \textbf{augmented matrix} \([A\mid b]\) records a linear system.
The three \textbf{elementary row operations} are
\[
\begin{array}{ll}
R_i\leftrightarrow R_j &(i\ne j),\\
R_i\leftarrow cR_i &(c\ne0),\\
R_i\leftarrow R_i+cR_j &(i\ne j).
\end{array}
\]
Each operation acts on an entire row, including the right-hand side.
\notetarget{note:01:row-equivalence}Two matrices are \textbf{row equivalent} if a finite sequence of these
operations changes one into the other; the empty sequence is allowed.

\begin{proposition}[Preservation of the solution set]\label{prop:01:row-operations}\notetarget{note:01:row-preservation}
An elementary row operation on \([A\mid b]\) preserves exactly the
solutions of \(Ax=b\).
\end{proposition}
\begin{proof}\notetarget{note:01:row-preservation-proof}
Interchanging equations changes only their order. Multiplying an
equation by nonzero \(c\) is undone by multiplying by \(1/c\).
For row replacement the old equations
\[
\sum_k a_{ik}x_k=b_i,\qquad \sum_k a_{jk}x_k=b_j
\]
imply
\[
\sum_k(a_{ik}+ca_{jk})x_k=b_i+cb_j,\qquad
\sum_k a_{jk}x_k=b_j.
\]
Subtracting \(c\) times the unchanged second equation from the new first
recovers the old first equation. Thus the implication holds in both
directions for each operation, and hence for any finite sequence.
\end{proof}
\notetarget{note:01:row-operation-cautions}Multiplication of a row by zero is excluded: it could erase an equation.
In a replacement, the row being added remains unchanged.

\subsection*{Echelon form}
\notetarget{note:01:echelon-definition}A matrix is in \textbf{row echelon form} if all zero rows are at the
bottom and the first nonzero entry of each successive nonzero row occurs
strictly farther to the right. These first nonzero entries are the
\textbf{pivots}; they need not be \(1\).

\begin{theorem}[Finite elimination]\label{thm:01:echelon}\notetarget{note:01:echelon}
Every real matrix is row equivalent to a matrix in row echelon form.
\end{theorem}
\begin{proof}\notetarget{note:01:echelon-proof}
\begin{itemize}
\item If the matrix is zero, stop. Otherwise choose its leftmost nonzero
column \(j\). Swap rows to put a nonzero entry \(p\) in position \((1,j)\).
All earlier columns are zero.
\item For each \(i>1\), replace \(R_i\) by \(R_i-(a_{ij}/p)R_1\).
The new entry in column \(j\) is \(a_{ij}-(a_{ij}/p)p=0\); for \(k>j\),
\[
a'_{ik}=a_{ik}-\frac{a_{ij}}p a_{1k}\quad(k>j).
\]
\notetarget{note:01:elimination-shape}The resulting shape is
\[
\begin{pmatrix}
0&\cdots&0&p&* &\cdots&*\\
0&\cdots&0&0&&&\\
\vdots&&\vdots&\vdots&&B&\\
0&\cdots&0&0&&&
\end{pmatrix}.
\]
\item \notetarget{note:01:elimination-termination}Repeat within the lower-right matrix \(B\).
Operations between its rows preserve their earlier zeros, so each new
pivot is strictly to the right of the preceding pivot.
Each step fixes one row and removes at least one column from the remaining
problem. After at most \(\min(m,n)\) pivot steps, the remaining matrix
is zero or has no rows or columns. The nonzero rows now have increasing
pivot columns, and any remaining rows are zero: this is echelon form.
\end{itemize}
\end{proof}

\begin{example}[A skipped column and a redundant row]\label{example:01:echelon}\notetarget{note:01:echelon-example}
\notetarget{note:01:echelon-example-solution}Reduce the following matrix, keeping the row operations visible:
\begin{align*}
\begin{pmatrix}0&0&1&2&3\\2&4&-2&0&4\\2&4&-1&2&7\end{pmatrix}
&\xrightarrow{R_1\leftrightarrow R_2}
\begin{pmatrix}2&4&-2&0&4\\0&0&1&2&3\\2&4&-1&2&7\end{pmatrix}\\
&\xrightarrow{R_3\leftarrow R_3-R_1}
\begin{pmatrix}2&4&-2&0&4\\0&0&1&2&3\\0&0&1&2&3\end{pmatrix}\\
&\xrightarrow{R_3\leftarrow R_3-R_2}
\begin{pmatrix}2&4&-2&0&4\\0&0&1&2&3\\0&0&0&0&0\end{pmatrix}\\
&\xrightarrow{R_1\leftarrow R_1/2}
\begin{pmatrix}1&2&-1&0&2\\0&0&1&2&3\\0&0&0&0&0\end{pmatrix}.
\end{align*}
\notetarget{note:01:echelon-observations}The pivot columns are 1 and 3. Column 2 is not zero, but contains no
pivot. A further replacement \(R_1\leftarrow R_1+R_2\) clears the entry
above the second pivot; this extra step is not required for echelon form.
\end{example}

\notetarget{note:01:pivot-variables}For \(Ax=b\), eliminate the coefficient columns and carry \(b\) along.
Write the result as \([U\mid d]\), where \(U\) is in echelon form.
Suppose its nonzero rows have pivot columns \(j_1<\cdots<j_r\).
The unknowns \(x_{j_i}\) are the \textbf{pivot variables}; the remaining
unknowns are \textbf{free variables}.

\begin{theorem}[Consistency and free variables]\label{thm:01:consistency}\notetarget{note:01:consistency}
The system \(Ax=b\) is consistent if and only if every zero row of
\(U\) has corresponding right-hand entry zero. If it is consistent,
every choice of the \(n-r\) free variables gives exactly one solution,
and every solution is obtained this way. A consistent system has a unique
solution exactly when there are no free variables.
\end{theorem}
\begin{proof}\notetarget{note:01:consistency-proof}
\begin{itemize}
\item A zero coefficient row requires \(0=d_i\), proving necessity.
\item \notetarget{note:01:back-substitution}Suppose all these conditions hold. Assign arbitrary real values to the
free variables. For \(i=r,r-1,\ldots,1\), set
\[
\boxed{x_{j_i}=\frac{d_i-\sum_{k>j_i}u_{ik}x_k}{u_{i j_i}}.}
\]
The pivot \(u_{i j_i}\) is nonzero. Every variable on the right is free
or has already been determined in a lower row. Each step therefore gives
one value and satisfies row \(i\); the zero rows already hold.
Reversibility gives a solution of \(Ax=b\).
\item \notetarget{note:01:all-solutions-proof}Any solution has specific free-coordinate values. The same formulas
force its pivot coordinates, so every solution is obtained.
Different free-coordinate choices give different solutions.
There is exactly one solution when there are no free variables;
otherwise varying a free coordinate gives infinitely many.
\end{itemize}
\end{proof}

\begin{example}[All coefficients in a representation]\label{example:01:all-coefficients}\notetarget{note:01:all-coefficients}
\notetarget{note:01:all-coefficients-solution}For \(A,b\) in Example~\ref{example:01:representation},
\[
\left[\begin{array}{rrrr|r}2&-4&0&0&4\\-1&2&0&0&-2\\0&0&1&2&-3\end{array}\right]
\xrightarrow{R_2\leftarrow R_2+R_1/2}
\left[\begin{array}{rrrr|r}2&-4&0&0&4\\0&0&0&0&0\\0&0&1&2&-3\end{array}\right].
\]
Interchange rows 2 and 3. The nonzero equations are
\[
2x_1-4x_2=4,\qquad x_3+2x_4=-3.
\]
\notetarget{note:01:all-coefficients-family}The free variables are \(x_2=s\), \(x_4=t\), giving all solutions
\[
\boxed{x=\begin{pmatrix}2+2s\\s\\-3-2t\\t\end{pmatrix}
=\begin{pmatrix}2\\0\\-3\\0\end{pmatrix}
+s\begin{pmatrix}2\\1\\0\\0\end{pmatrix}
+t\begin{pmatrix}0\\0\\-2\\1\end{pmatrix},\quad s,t\in\R.}
\]
\notetarget{note:01:all-coefficients-check}Substitution gives \(2(2+2s)-4s=4\),
\(-(2+2s)+2s=-2\), and \((-3-2t)+2t=-3\).
Every solution has this form because its second and fourth coordinates
are precisely \(s,t\).

\notetarget{note:01:inconsistent-right-hand-side}For right-hand side \(c=(1,0,0)^\top\), the same first row operation
instead produces a zero coefficient row with right-hand entry \(1/2\).
It requires \(0=1/2\), proving inconsistency.
\end{example}

\begin{proposition}[More unknowns than equations]\label{prop:01:excess}\notetarget{note:01:excess}
If \(A\in\R^{m\times n}\) and \(n>m\), then \(Ax=0\) has a nonzero solution.
\end{proposition}
\begin{proof}\notetarget{note:01:excess-proof}
Row operations leave a zero right-hand side zero, so the homogeneous
system is consistent. There are at most \(m\) pivots and hence fewer
than \(n\). Set one free variable to \(1\), the other free variables
to \(0\), and obtain the pivot variables by back substitution.
This gives a solution whose chosen free coordinate is nonzero.
\end{proof}
\notetarget{note:01:excess-not-necessary}The condition is sufficient, not necessary. For example,
\(x_1+x_2=0\), \(2x_1+2x_2=0\) have solutions \((t,-t)^\top\)
although there are two equations and two unknowns; substitution and the
first equation prove that this is their entire solution set.

\section*{4. Column spaces, null spaces, and translated solutions}

\notetarget{note:01:spaces}For \(A\in\R^{m\times n}\), define the \textbf{column space} and
\textbf{null space} by
\[
\mathcal C(A)=\{Ax:x\in\R^n\}\subseteq\R^m,\qquad
N(A)=\{x\in\R^n:Ax=0\}\subseteq\R^n.
\]
Membership \(b\in\mathcal C(A)\) means that \(b\) is a combination of
\(A\)'s columns, equivalently that \(Ax=b\) is consistent.
An element of \(N(A)\) lists coefficients of a zero combination of those
columns. The two sets generally lie in different coordinate spaces.

\begin{proposition}[Closure under linear combinations]\label{prop:01:closure}\notetarget{note:01:closure}
The sets \(\mathcal C(A)\) and \(N(A)\) contain zero and are closed under
linear combinations.
\end{proposition}
\begin{proof}\notetarget{note:01:closure-proof}
The equality \(A0=0\) puts zero in both sets. If \(u=Ax\), \(v=Ay\), then
\[
su+tv=sAx+tAy=A(sx+ty)\in\mathcal C(A).
\]
If \(Ax=Ay=0\), then \(A(sx+ty)=sAx+tAy=0\), so
\(sx+ty\in N(A)\). \notetarget{note:01:finite-closure}Repeating the two-term calculation proves closure
for any finite combination; the empty one is zero.
\end{proof}

\begin{theorem}[The whole solution set]\label{thm:01:translate}\notetarget{note:01:translate}
Suppose \(Ax_0=b\). Then
\[
\{x:Ax=b\}=x_0+N(A):=\{x_0+z:z\in N(A)\}.
\]
\end{theorem}
\begin{proof}\notetarget{note:01:translate-proof}
If \(Ax=b\), then \(A(x-x_0)=Ax-Ax_0=b-b=0\). Hence
\(x=x_0+z\) with \(z=x-x_0\in N(A)\).
Conversely, if \(Az=0\), then \(A(x_0+z)=Ax_0+Az=b\).
These two implications prove equality of the sets.
\end{proof}

\begin{example}[The sets for a rectangular matrix]\label{example:01:spaces}\notetarget{note:01:rectangular-spaces}
For the matrix of Example~\ref{example:01:representation},
\[
\mathcal C(A)=
\left\{\begin{pmatrix}2u\\-u\\v\end{pmatrix}:u,v\in\R\right\}
=\left\{b\in\R^3:b_1=-2b_2\right\}.
\]
\notetarget{note:01:rectangular-column-space-proof}Indeed \(Ax=(2u,-u,v)^\top\) with \(u=x_1-2x_2\),
\(v=x_3+2x_4\). Conversely every \(u,v\) is obtained by taking
\(x=(u,0,v,0)^\top\). The displayed equation on \(b\) is equivalent
to this form by setting \(u=-b_2\), \(v=b_3\).
\notetarget{note:01:rectangular-null-space}Also,
\[
N(A)=\left\{
 s\begin{pmatrix}2\\1\\0\\0\end{pmatrix}
+t\begin{pmatrix}0\\0\\-2\\1\end{pmatrix}:s,t\in\R
\right\},
\]
because \(Ax=0\) says \(x_1=2x_2\), \(x_3=-2x_4\).
Thus \(\mathcal C(A)\subseteq\R^3\), whereas \(N(A)\subseteq\R^4\).
\notetarget{note:01:rectangular-translated-solutions}For \(b=(4,-2,-3)^\top\), adding the particular solution
\((2,0,-3,0)^\top\) to this null space reproduces all the coefficients
in Example~\ref{example:01:all-coefficients}.
\end{example}

\section*{5. Independence and unique representation}

\notetarget{note:01:independence}A finite list \(v_1,\ldots,v_k\) of columns in the same \(\R^m\) is
\textbf{linearly independent} if
\[
c_1v_1+\cdots+c_kv_k=0\quad\Longrightarrow\quad c_1=\cdots=c_k=0.
\]
It is \textbf{linearly dependent} if a zero combination exists with at
least one coefficient nonzero. The same definition applies to any list
of matrices of the same size. The empty list is independent.

\notetarget{note:01:standard-independence}The standard columns \(e_1,\ldots,e_m\) are independent: coordinate
\(i\) of \(\sum_jc_je_j=0\) says \(c_i=0\). \notetarget{note:01:zero-member}A list containing zero
is dependent by giving that member coefficient \(1\) and all other
members coefficient \(0\). \notetarget{note:01:independent-sublist}A sublist of an independent list is
independent: extend any relation in the sublist by zero coefficients
on the omitted members and apply the original independence.

\begin{proposition}[Dependence and a redundant member]\label{prop:01:redundancy}\notetarget{note:01:redundancy}
A nonempty list is dependent if and only if some member is a linear
combination of the others.
\end{proposition}
\begin{proof}\notetarget{note:01:redundancy-proof}
In a nonzero relation \(\sum_i c_iv_i=0\), choose \(j\) with \(c_j\ne0\).
Solving for that member gives
\[
v_j=-\sum_{i\ne j}\frac{c_i}{c_j}v_i.
\]
Conversely, from \(v_j=\sum_{i\ne j}d_iv_i\) obtain a zero combination
with coefficient \(1\) on \(v_j\) and \(-d_i\) on the others.
For a one-member list the combination of the others is the empty sum,
so both statements mean \(v_1=0\).
\end{proof}

\needspace{10\baselineskip}
\begin{theorem}[Independence and uniqueness]\label{thm:01:independence}\notetarget{note:01:independence-equivalence}
Let \(A=[\,a_1\ \cdots\ a_n\,]\). The following conditions are equivalent:
\begin{enumerate}
\item The columns \(a_1,\ldots,a_n\) are linearly independent.
\item \(N(A)=\{0\}\).
\item For every \(b\in\mathcal C(A)\), the system \(Ax=b\) has exactly one solution.
\end{enumerate}
\end{theorem}
\begin{proof}\notetarget{note:01:independence-equivalence-proof}
\begin{itemize}
\item (a) \(\Longleftrightarrow\) (b): \(Ax=\sum_jx_ja_j\), so both conditions
say that \(Ax=0\) forces \(x=0\).
\item (b) \(\Longrightarrow\) (c): for \(b\in\mathcal C(A)\), choose \(x_0\)
with \(Ax_0=b\). Theorem~\ref{thm:01:translate} gives the solution set
\(x_0+N(A)=x_0+\{0\}=\{x_0\}\).
\item (c) \(\Longrightarrow\) (b): take \(b=0\in\mathcal C(A)\).
Since \(A0=0\) and the solution is unique, \(N(A)=\{0\}\).
\end{itemize}
\end{proof}
\notetarget{note:01:attainability-restriction}The restriction \(b\in\mathcal C(A)\) is essential. For
\(A=(1,0)^\top\), the equation \(Ax=0\) forces \(x=0\), but
\(Ax=(0,1)^\top\) is impossible by its second coordinate.

\notetarget{note:01:too-many-columns}A list of more than \(m\) columns in \(\R^m\) is dependent:
put the columns into an \(m\times n\) matrix with \(n>m\) and apply
Proposition~\ref{prop:01:excess} to obtain a nonzero coefficient relation.

\begin{example}[Dependence at exceptional parameter values]\label{example:01:parameter}\notetarget{note:01:parameter}
For which real \(k\) are
\[
v_1=\begin{pmatrix}k\\1\\0\end{pmatrix},\qquad
v_2=\begin{pmatrix}1\\k\\1\end{pmatrix},\qquad
v_3=\begin{pmatrix}0\\1\\k\end{pmatrix}
\]
independent? \notetarget{note:01:parameter-solution}A relation \(\alpha v_1+\beta v_2+\gamma v_3=0\) says
\[
k\alpha+\beta=0,\qquad \alpha+k\beta+\gamma=0,\qquad
\beta+k\gamma=0.
\]
\notetarget{note:01:parameter-elimination}The first two equations give, without dividing by \(k\),
\[
\beta=-k\alpha,\qquad \gamma=(k^2-1)\alpha.
\]
Substitution in the third gives
\[
0=-k\alpha+k(k^2-1)\alpha=k(k^2-2)\alpha.
\]
\notetarget{note:01:parameter-classification}If \(k\notin\{0,\sqrt2,-\sqrt2\}\), this forces \(\alpha=0\), and then
\(\beta=\gamma=0\). The vectors are independent.
At each exceptional value, take \(\alpha=1\); the nonzero relation is
\[
v_1-kv_2+(k^2-1)v_3=0.
\]
Its three coordinates are \(k-k=0\),
\(1-k^2+k^2-1=0\), and \(k(k^2-2)=0\), respectively.
Thus dependence occurs exactly for \(k=0,\pm\sqrt2\).
\end{example}
